\documentclass[12pt, a4paper]{article}

\usepackage{amsmath, amssymb, amsfonts}
\usepackage{float}
\usepackage{graphicx}
\usepackage{listings}
\usepackage{xcolor}
\usepackage[linesnumbered, ruled]{algorithm2e}
\usepackage{subcaption}
\usepackage[hidelinks]{hyperref}
% Each algorithm is written twice: the algorithm2e block (what the PDF typesets)
% and a plain-text Verbatim block right after it (what the web shows — pandoc
% renders Verbatim as a copyable <pre> code block). \excludecomment makes pdflatex
% SKIP the Verbatim blocks, so the PDF shows only the algorithm2e version; pandoc
% conversely drops each algorithm2e to an (empty, CSS-hidden) div and keeps the
% Verbatim. Because pdflatex never sees the Verbatim, its Unicode maths glyphs
% (⟨⟩ ← τ ∈ ∞ …) need no inputenc/literate setup.
\usepackage{comment}
\excludecomment{Verbatim}

\definecolor{codegreen}{rgb}{0,0.6,0}
\definecolor{codegray}{rgb}{0.5,0.5,0.5}
\definecolor{codepurple}{rgb}{0.58,0,0.82}
\hypersetup{linktoc=all}

\lstset{
    language=Python,
    basicstyle=\ttfamily\small,
    keywordstyle=\color{blue},
    commentstyle=\color{codegreen},
    stringstyle=\color{codepurple},
    breaklines=true,
    numbers=left,
    numberstyle=\tiny\color{codegray},
    frame=single,
    showstringspaces=false
}

\title{AI Masterclass Test 1}
\author{}
\date{}


\begin{document}

\section{Problem Formulation}

\noindent A problem, $\mathcal{P}$, is defined by the following elements:

\[
\mathcal{P} = \langle S, s_0, A, \tau, c, G \rangle
\]

\noindent where:
\begin{itemize}
    \item $S$ is the state space
    \item $s_0 \in S$ is the initial state
    \item $A = \{a_1, \dots, a_k\}$ is set of actions or operations
    \item $\tau : S \times A \to S$ is the state transition function
    \item $c : S \times A \times S \to \mathbb{R}$ is the cost function
    \item $G \subseteq S$ defines the goal to be achieved.
\end{itemize}

\noindent A solution 
\[
\vec{a} = (a_0, a_1, \dots, a_n) \in \text{soln}(\mathcal{P})
\]
which induces a state sequence
\[
(s_0, s_1, \dots, s_{n+1})
\]
has an associated cost:
\[
c(\vec{a}) = \sum_{i=0}^{n} c(s_i, a_i, s_{i+1})
\]

\noindent A solution $\vec{a}^*$ is \textbf{optimal} if it minimises corresponding cost:
\[
\vec{a}^* \in \arg \min_{\vec{a} \in \text{soln}(\mathcal{P})} c(\vec{a})
\]

\end{document}
